\section{Introduction}

\label{sec:introduction}

Modern Bayesian inference often requires efficient sampling from complex, high-dimensional, and multimodal posterior distributions. As these problems grow in scale and complexity, traditional sampling methods face significant computational and statistical challenges. The advent of automatic differentiation has enabled the widespread use of gradient-based samplers, making methods such as stochastic gradient Langevin dynamics (SGLD) and stochastic gradient Hamiltonian Monte Carlo (SGHMC) practical and popular choices for large-scale applications \cite{welling2011sgld,chen2014sghmc}. While these approaches advance a single Markov chain using noisy gradients, an alternative paradigm is offered by interacting particle methods, which evolve a population of particles whose empirical distribution approximates the target posterior. This population-based view is central to particle-based variational inference \cite{pmlr-v97-liu19i}. Within this class, SVGD is a prominent deterministic representative, and recent work has examined acceleration, particle degeneracy, and kernel design in related Stein-particle methods \cite{pmlr-v97-liu19i,pmlr-v80-zhuo18a,wang2019matrixsvgd}. It opens new avenues for improving sampling efficiency and robustness, especially in challenging inference settings.

A useful unifying perspective is to view interacting particle methods as
discretizations of a particle flow (possibly with a stochastic component) that
transports an initial ensemble toward the target distribution \cite{Sandu_2025_EnsembleFokkerPlanck}.
Here, we focus on Stein variational gradient descent (SVGD) \cite{liu2016svgd} as a
representative interacting particle method. SVGD constructs a velocity field
for the particles based on the target density and a reproducing kernel Hilbert
space (RKHS) structure, and updates the particles by discretizing this flow \cite{liu2016svgd}.
%The explicit form of the SVGD update and its velocity field are given in
%\cref{eq:svgd-update,eq:svgd-velocity} in Section~\ref{sec:method}. 
The velocity field naturally separates into an attractive drift, which moves
particles toward high-density regions, and a repulsive interaction, which
maintains particle diversity and support coverage of multiple modes.

This drift-repulsion structure creates the main numerical challenge studied in
this paper. When particles crowd, repulsion can become stiff and destabilize
updates. On ther hand, the drift can require finer resolution on anisotropic targets to
follow the geometry of $\log p$ closely. A single step size must therefore
balance competing requirements and can either over-damp repulsion, leading to
loss of diversity and mode coverage, or under-resolve drift, leading to slow
progress.

Moreover, repulsion typically requires kernel evaluation at each
step, so increasing the stability by reducing  the step-size can become
computationally expensive. Finally, repulsive forces can weaken in high
dimensions or under limited particle resolution, leading to particle collapse
and poor support of multiple modes \cite{pmlr-v80-zhuo18a,pmlr-v108-zhang20d}.
In practice, the resulting single-step-size compromise is both statistically
and computationally inefficient: it can degrade support coverage and predictive
quality while still spending substantial kernel work to stabilize the wrong
part of the update.

%This paper focuses specifically on this numerical weakness of SVGD: a single global
%step size must simultaneously resolve attractive drift and
%diversity-preserving repulsion, even though these components can evolve on
%different effective time scales.

A common approach to solving systems with significantly different 
components is to use partitioned methods such as implicit-explicit (IMEX)  \cite{Pareschi_2000,sarshar2020parimex} and
multirate integration methods \cite{Sandu_2016_MR-GARK}.  In the context of ordinary
differential equations (ODEs), multirate methods exploit a natural splitting of the dynamics into 
fast and slow components, allowing each to be integrated with a step size appropriate
to its scale \cite{sarshar2019mrgark,sarshar2021implicitgark,sarshar2021glmadi}. 
For example, in a split ODE of the form $\dot{x} = \phi_{\text{fast}}(x) + \phi_{\text{slow}}(x)$, multirate schemes
apply more frequent, smaller steps to the stiff or rapidly varying part, while updating the
slow part less often. This approach improves both stability and efficiency in multi-scale
dynamical systems, and has been successfully applied in a variety of scientific computing contexts \cite{sarshar2020parimex,sarshar2022surrogate}. 

Motivated by advances in multirate time-integration
for multi-scale dynamical systems—including MrGARK, IMEX, and related splitting schemes \cite{sarshar2019mrgark,sarshar2021implicitgark,sarshar2021glmadi} we propose to integrate the drift and repulsion terms on separate time scales. This leads to a number of multirate SVGD variants, including symmetric splitting, fixed repulsion substepping, and adaptive error-controlled drift substepping.

To assess the effectiveness of these methods, we conduct a comprehensive empirical evaluation. We compare our multirate SVGD variants with established approaches from the literature across a range of benchmark problems, using a variety of task-appropriate metrics. Results highlight the strengths and limitations of each method under diverse geometric and statistical challenges.

\subsection*{Contributions}

The main contributions of this work are:
\begin{itemize}
  \item Multirate SVGD formulations that explicitly separate drift and repulsion
  and support symmetric splitting and fixed multirate schedules.
  \item An adaptive error-controlled multirate SVGD variant that selects 
  substeps based on a local error estimator while keeping repulsion stable.
  \item A broad empirical study showing that multirate SVGD improves robustness
  and quality-cost tradeoffs compared to SVGD and representative chain methods, with the strongest gains on
  anisotropic, multimodal, and hierarchical targets.
  \item A benchmark suite spanning many different targets,
  equipped with appropriate diagnostics, early stopping, performance evaluation metrics for efficiency comparisons.
\end{itemize}

We present our work as follows. Section~\ref{sec:method} presents the SVGD
variants and the multirate construction. Section~\ref{sec:exp_setup} describes
the benchmarks, diagnostics, and experimental protocol. Section~\ref{sec:results}
reports the results. Section~\ref{sec:conclusions} summarizes the findings and
discusses directions for future work.
